
% !TEX program = xelatex
% Overleaf 编译设置: Menu -> Compiler -> XeLaTeX
% 不要用 pdfLaTeX, 否则 ctex 会报 "fontset fandol unavailable"。
\documentclass[10pt,twocolumn,fontset=none]{ctexart}
\usepackage{iftex}
\ifPDFTeX
  \errmessage{This document must be compiled with XeLaTeX. In Overleaf: Menu -> Compiler -> XeLaTeX}
\fi
\usepackage[a4paper,top=1.75cm,bottom=1.75cm,left=1.55cm,right=1.55cm,columnsep=0.72cm]{geometry}
\usepackage{fontspec}
\usepackage{xeCJK}
% 字体兼容: 优先使用本地 Windows SimSun; 其次使用 Overleaf 常见 Noto CJK; 若没有则回退到 Fandol 或其他。
\IfFontExistsTF{SimSun}{
  \setCJKmainfont{SimSun}
}{
  \IfFontExistsTF{Noto Serif CJK SC}{
    \setCJKmainfont{Noto Serif CJK SC}
  }{
    \IfFontExistsTF{FandolSong-Regular}{
      \setCJKmainfont{FandolSong-Regular}
    }{
      \setCJKmainfont{AR PL UMing CN}
    }
  }
}
\IfFontExistsTF{SimHei}{
  \setCJKsansfont{SimHei}
}{
  \IfFontExistsTF{Noto Sans CJK SC}{
    \setCJKsansfont{Noto Sans CJK SC}
  }{
    \IfFontExistsTF{FandolHei-Regular}{
      \setCJKsansfont{FandolHei-Regular}
    }{
      \setCJKsansfont{AR PL UKai CN}
    }
  }
}
\usepackage{amsmath,amssymb,bm}
\usepackage{graphicx}
\usepackage{booktabs}
\usepackage{multirow}
\usepackage{array}
\usepackage{enumitem}
\usepackage{tikz}
\usetikzlibrary{arrows.meta,positioning,fit,calc,shapes.geometric,shapes.misc,backgrounds}
\usepackage[hidelinks]{hyperref}
\usepackage{balance}
\usepackage{caption}
\usepackage{url}
\captionsetup{font=small,labelfont=bf}
\setlength{\parskip}{0pt}
\setlength{\parindent}{1.5em}
\setlist[itemize]{leftmargin=1.4em,itemsep=0pt,topsep=2pt}
\setlist[enumerate]{leftmargin=1.4em,itemsep=0pt,topsep=2pt}
\newcommand{\method}{SEER-POI}
\newcommand{\methodfull}{Spatio-functional Evidence-routed Expert Retrieval-augmented framework for Next Point-of-Interest Recommendation}
\newcommand{\hsid}{HSID}
\newcommand{\blank}{\rule{1.15cm}{0.35pt}}
\newcommand{\wideblank}{\rule{1.8cm}{0.35pt}}
\newcommand{\code}[1]{\texttt{#1}}
\title{\vspace{-1.0em}\bfseries \method: 基于时空与功能语义证据路由检索增强的下一兴趣点预测框架}
\author{匿名作者}
\date{}
\begin{document}
\maketitle
\vspace{-1.35em}

\begin{abstract}
下一兴趣点(Point-of-Interest, POI)预测试图根据用户历史轨迹推断其下一次最可能访问的位置。该任务表面上可以被写成一个序列分类问题, 但真实城市移动行为远比“历史 ID 到未来 ID”的映射更复杂。一个 POI 既是独立地点, 又属于某个区域、类别、商圈与功能语义结构; 一次移动既可能由地理可达性驱动, 也可能由惯常转移、时间偏好或语义探索驱动。现有下一 POI 方法虽然在序列建模、时空图建模和大语言模型推理方面取得进展, 但仍存在两个紧密耦合的瓶颈: 第一, POI 常被建模为平坦离散 ID, 难以显式表达区域层次、类别功能和语义邻域之间的共享关系; 第二, 候选生成常依赖单一路径检索或静态融合, 在稀疏轨迹和长尾 POI 场景下容易漏召回真实目标, 从而限制后续排序器的性能上限。

本文提出 \method, 一个面向下一 POI 预测的层次语义证据路由候选生成框架。与直接以“语义 ID”或“RAG 专家”命名不同, \method 将核心贡献重新组织为一条完整故事线: 先用层次语义标识(\hsid)把孤立 POI 转化为可共享的空间-功能语义节点, 再用多证据路由候选生成模块根据当前轨迹上下文动态融合空间、语义、转移和时间四类证据, 最后在去重、池化和压缩后的候选集合上执行候选条件化排序。该框架强调候选生成与最终预测之间的功能边界: 检索阶段负责提高真实目标覆盖率并压缩搜索空间, 排序阶段负责在高相关候选内部进行细粒度判别。本文给出完整的问题定义、方法建模、框架图、算法流程、复杂度分析以及 CCF-A 会议论文风格的实验表格模板。实验数值位置均按要求留空, 以便后续填入正式结果。
\end{abstract}

\noindent\textbf{关键词:} 下一兴趣点预测; 层次语义标识; 证据路由; 候选生成; 检索增强推荐; 轨迹建模

\vspace{0.35em}
\noindent\textbf{Abstract.} Next point-of-interest prediction is not merely an ID-to-ID sequence classification problem. A destination carries hierarchical spatial semantics, and a mobility decision may be supported by spatial accessibility, behavioral transition, temporal preference, or semantic exploration. This paper proposes \method, a Spatio-functional Evidence-routed Expert Retrieval-augmented (SEER) framework for next POI recommendation. \method{} first constructs compact hierarchical semantic identifiers (HSID) for POIs, then dynamically routes spatial, semantic, transition, and temporal evidence through dedicated expert channels to generate a compact, high-recall candidate set, and finally performs candidate-conditioned ranking.

\section{引言}

下一兴趣点预测是位置服务、个性化推荐、城市计算和移动行为理解中的基础任务。给定用户已经发生的签到轨迹, 系统需要判断用户下一次最可能访问哪一个 POI。这个问题看似直接: 只需要把历史序列输入模型, 让模型输出下一个 POI 的编号。但在真实城市中, “位置”并不是一个单纯的符号, “下一步”也不是一个只由序列顺序决定的标签。用户早晨从住宅区到地铁站再到办公区, 体现的是空间可达性和稳定通勤规律; 用户午间从办公楼转向附近餐饮区, 体现的是时段偏好和功能区域关系; 用户周末从常去商圈转向另一个风格相近的商圈, 体现的是语义相似与探索行为。若模型只记住 POI ID 的共现频率, 就很难解释这些行为背后的多源证据。

近年来, 下一 POI 预测经历了从浅层统计模型到深度序列模型, 再到图神经网络、对比学习和大语言模型的演进。循环神经网络与 Transformer 提升了时序依赖建模能力; 图模型强调用户、POI 与类别之间的结构关系; 大语言模型进一步带来了自然语言化轨迹描述、候选解释和结构化输出能力。与此同时, 检索增强范式开始被用于缩小候选空间, 使大模型不再需要在全体 POI 上自由生成答案。然而, 我们在代码实现和误差诊断中发现, 仅仅换用更强排序器或更长提示词并不能根本解决下一 POI 任务中的核心困难。真正决定系统上限的, 往往是两个更底层的问题。

第一个问题是 POI 表示过于平坦。传统系统通常把每个 POI 当作彼此独立的离散 ID, 最多再附加类别、经纬度或文本描述。这种表示方式能区分地点个体, 但难以表达地点之间的层次共享关系。例如, 两个 POI 可能位于同一粗区域、相邻细区域、同一功能类别或同一语义簇。对于高频 POI, 模型可以从重复访问中学习表示; 对于低频或长尾 POI, 如果缺乏显式共享结构, 模型就很难借助相似 POI 的统计信号完成泛化。换言之, 平坦 ID 让模型知道“这是哪个点”, 却没有稳定告诉模型“这个点处于怎样的空间和功能结构中”。

第二个问题是候选生成缺乏上下文自适应性。真实城市候选空间通常很大, 若直接让模型在全体 POI 中预测, 计算代价和噪声都会很高。因此实际系统往往先进行候选召回, 再执行精排。但候选召回本身并不简单。语义向量检索可能召回功能相似但距离很远的 POI; 空间近邻可能召回地理接近但行为上无关的 POI; 历史频次候选可能偏向用户常去地点而忽略探索访问; 时间偏好候选又可能缺乏具体空间约束。更关键的是, 如果真实目标没有进入候选池, 后续排序器无论多强都无法恢复正确答案。因此, 候选生成不是排序前的工程细节, 而是下一 POI 预测的核心建模环节。

上述两个问题彼此耦合。若 POI 没有层次语义表示, 检索阶段就难以识别语义相关但观测稀少的位置; 若候选池质量不足, 层次语义信息也无法在最终排序中发挥作用。因此, 本文不把下一 POI 预测理解为单纯“做一个更强的模型”, 而是从系统链路角度重新组织任务: 首先让 POI 具有可共享、可解释且可压缩的层次语义表示; 然后让候选生成根据当前轨迹在多类证据之间动态选择; 最后让排序器只在高覆盖、高相关的候选集合中进行细粒度判断。

基于这一思路, 本文提出 \method。名称中的 SEER 表示 Spatio-functional Evidence-routed Expert Retrieval-augmented, 强调本文的两个核心关键词: 层次语义与多源证据路由。与已有直接使用 Semantic ID 作为生成目标的工作不同, 本文的层次语义标识不是让模型生成一串 ID token, 而是贯穿检索索引、候选表示、转移统计和排序提示构造的共享语义支架。与直接命名为 ExpertRAG 的做法不同, 本文避免把贡献表述为通用 RAG 名称, 而是将其落到下一 POI 场景中的候选生成机制: 空间、语义、转移和时间四类证据通道如何在不同轨迹上下文中被动态融合。

本文贡献如下。
\begin{itemize}
    \item 提出面向下一 POI 预测的层次语义标识 \hsid, 将 POI 从孤立离散编号扩展为由粗区域、细区域、功能类别与语义簇组成的层次语义节点, 为长尾 POI 和稀疏轨迹提供可共享的结构信号。
    \item 提出多证据路由候选生成机制, 将候选召回分解为空间可达性、语义相似性、行为转移规律和时间偏好四类证据通道, 并根据当前轨迹上下文动态融合, 缓解单一路径检索中的噪声挤占和真实目标漏召回问题。
    \item 设计候选去重、元数据池化和紧凑结构化提示方法, 使 \hsid 在进入大模型排序时保持语义增益, 同时避免候选描述膨胀导致的上下文截断和注意力稀释。
    \item 给出完整的 CCF-A 论文结构, 包括框架图、算法流程、公式化建模、复杂度分析、总体性能表、候选召回表、消融表和效率表, 便于后续直接补充正式实验结果。
\end{itemize}

\section{相关工作}

\subsection{下一 POI 预测}
下一 POI 预测长期以来是时空推荐的代表性任务。早期方法多基于马尔可夫链、矩阵分解或循环网络, 重点刻画用户连续签到中的局部转移规律。随着深度学习发展, 注意力机制和 Transformer 被引入轨迹序列建模, 用于捕获长程依赖、周期行为和上下文交互。图神经网络进一步把用户、POI、类别、区域和转移关系组织为多视图图结构, 能够挖掘更复杂的时空关联。近年来, 对比学习、预训练和生成式推荐方法被用于提升稀疏轨迹下的表示鲁棒性。

这些方法为下一 POI 预测提供了强大的轨迹编码基础, 但多数工作仍以平坦 POI ID 为核心预测单元。即便部分方法使用类别、距离或时间特征, 这些信息也常作为附加特征输入模型, 而不是形成可复用、可检索、可压缩的层次语义标识。本文与这些工作不同, 并不只增强序列编码器, 而是从候选生成链路前端重构 POI 表示, 使后续检索和排序都能利用统一的层次语义支架。

\subsection{语义 ID 与地理语义表示}
推荐系统中的语义 ID 试图缓解离散 ID 的稀疏性和不可解释性。GNPR-SID 提出基于 Semantic ID 的生成式下一 POI 推荐, 通过语义 POI ID 增强 LLM 对 POI 的理解能力\cite{gnprsid}. ROS 进一步提出 Hierarchical Spatial Semantic ID, 将 coarse-to-fine locality 与 POI semantics 离散为组合 token, 以增强 LLM 的地理推理能力\cite{ros}. 这些工作说明, 将 POI 从原始编号映射到语义结构中已经成为 next POI 研究中的重要方向。

本文保留层次语义表示的思想, 但避免把创新点直接命名为新的 SID。原因在于, 本文关注的不是“让模型生成语义 ID”, 而是“让层次语义作为证据贯穿候选生成和排序”。\hsid 既进入 POI 向量索引文本, 也进入候选元数据, 还可支持类别、区域和语义簇级别的多尺度转移统计。它更像是一种轻量的空间-功能语义骨架, 而不是最终输出空间本身。因此, 本文将方法命名为 \method, 以突出层次证据路由和候选生成机制。

\subsection{检索增强与大语言模型推荐}
大语言模型为下一 POI 预测带来了新的建模方式。模型可以读取自然语言化轨迹、候选 POI 表和约束格式, 输出可解释的下一位置预测。但 LLM 若缺少外部检索或地理约束, 容易生成泛化但空间上不合理的答案。RALLM-POI 将历史轨迹检索与地理距离重排结合, 用于零样本 next POI 推荐\cite{rallmpoi}. 这类工作说明检索增强对于 LLM-based POI 推荐是必要的。

然而, 现有检索增强方法往往更强调“检索一批相关上下文”, 对候选本身的多源证据结构讨论不足。本文认为, 对于下一 POI, 检索不仅是把历史样本拿给 LLM 参考, 更是构建最终候选决策空间。候选生成需要同时考虑地理、语义、转移和时间, 并根据当前轨迹动态选择证据。因此, 本文将 RAG 思想下沉为候选生成模块, 通过多证据通道和上下文路由形成更紧凑、更高覆盖的候选池。

\subsection{专家路由与多证据融合}
通用 RAG 领域已有工作使用 mixture-of-experts 或动态路由改进检索效率, 例如 ExpertRAG 将 MoE 与 RAG 结合, 讨论动态检索门控与专家路由\cite{expertrag}. 但该类工作主要面向知识密集型文本生成任务, 并不是针对下一 POI 预测中的候选池构造、地理约束、轨迹转移和 POI 层次语义而设计。

因此, 本文不沿用 ExpertRAG 作为方法名称, 而将对应思想重新表述为“多证据路由候选生成”。这一命名更贴近 POI 推荐任务本身: 多证据不是通用专家, 而是空间、语义、转移和时间四类明确的移动行为证据; 路由不是为了选择 LLM 内部专家, 而是为了在当前轨迹上下文下选择合适的候选召回来源。

\section{问题定义}

设用户集合为 $\mathcal U=\{u_1,u_2,\ldots,u_{|\mathcal U|}\}$, POI 集合为 $\mathcal P=\{p_1,p_2,\ldots,p_{|\mathcal P|}\}$。每个 POI $p\in\mathcal P$ 具有原始 ID $id_p$, 地理坐标 $g_p=(lat_p,lon_p)$, 类别 $c_p$, 以及可选属性 $a_p$, 如行政区域、子类别或文本描述。用户 $u$ 的第 $i$ 次签到记为
\begin{equation}
    x_{u,i}=\langle p_{u,i},t_{u,i}\rangle,
\end{equation}
其中 $p_{u,i}\in\mathcal P$ 为访问位置, $t_{u,i}$ 为时间戳。按时间排序后, 用户轨迹为
\begin{equation}
    T_u=\{x_{u,1},x_{u,2},\ldots,x_{u,n_u}\}.
\end{equation}
给定用户历史轨迹 $T^H_u$ 与近期上下文 $T^C_u$, 下一 POI 预测目标是预测真实下一位置 $p^*_u=p_{u,n_u}$。本文采用候选约束预测范式, 先生成候选集合 $C_u\subseteq\mathcal P$, 再在该集合上排序:
\begin{equation}
    \hat p_u=\arg\max_{p\in C_u} P(p\mid T^H_u,T^C_u,C_u).
\end{equation}
该形式显式区分候选生成与候选排序。候选生成阶段的目标是提高 $p^*_u\in C_u$ 的概率, 同时控制 $|C_u|$ 和噪声比例; 排序阶段的目标是在 $C_u$ 内将真实目标排到前列。若真实目标未进入候选池, 最终 Hit@K 必然为零。因此本文将候选召回率视为与最终排序指标同等重要的诊断指标。

为建模移动行为, 定义相邻签到之间的时间间隔和空间位移:
\begin{align}
    \Delta t_{u,i} &= t_{u,i}-t_{u,i-1},\\
    \Delta d_{u,i} &= \mathrm{Dist}(g_{p_{u,i-1}},g_{p_{u,i}}),
\end{align}
其中 $\mathrm{Dist}(\cdot,\cdot)$ 可采用 Haversine 距离。$\Delta t$ 与 $\Delta d$ 可用于判断样本更接近局部通勤、短距离消费、长距离探索还是稀疏不确定场景, 从而影响后续证据路由。

\section{方法}

\subsection{总体框架}
图~\ref{fig:framework} 展示了 \method 的整体框架。框架包含四个阶段。第一阶段离线构建 \hsid, 将每个 POI 映射为层次语义节点。第二阶段解析用户轨迹上下文, 提取最近位置、访问时间、类别序列、位移特征与重复访问特征。第三阶段执行多证据路由候选生成, 由空间、语义、转移和时间四个通道并行召回候选, 再根据上下文动态融合。第四阶段对候选池去重、池化元数据并构造紧凑候选表, 最后由候选条件化排序器输出下一 POI。

\begin{figure*}[t]
\centering
\resizebox{0.97\textwidth}{!}{
\begin{tikzpicture}[
    node distance=0.78cm and 0.98cm,
    box/.style={draw, rounded corners, align=center, minimum height=0.90cm, minimum width=2.45cm, font=\small},
    smallbox/.style={draw, rounded corners, align=center, minimum height=0.70cm, minimum width=2.10cm, font=\scriptsize},
    group/.style={draw, rounded corners, inner sep=0.25cm},
    arrow/.style={-{Latex[length=2.2mm]}, thick}
]
\node[box] (traj) {用户轨迹\\历史偏好 + 近期上下文};
\node[box, below=of traj] (poi) {POI 元数据\\坐标/类别/访问日志};

\node[group, right=1.0cm of traj, minimum width=3.15cm, minimum height=2.85cm, label={[font=\small]above:\textbf{层次语义标识}}] (hgroup) {};
\node[smallbox] (h1) at ($(hgroup.center)+(0,0.86)$) {粗区域};
\node[smallbox, below=0.14cm of h1] (h2) {细区域};
\node[smallbox, below=0.14cm of h2] (h3) {类别};
\node[smallbox, below=0.14cm of h3] (h4) {语义簇};

\node[box, right=0.95cm of hgroup] (query) {轨迹上下文解析\\查询表示 $q_u$};

\node[group, right=1.0cm of query, minimum width=5.65cm, minimum height=3.35cm, label={[font=\small]above:\textbf{多证据路由候选生成}}] (egroup) {};
\node[smallbox] (sp) at ($(egroup.center)+(-1.65,0.78)$) {空间证据\\距离/可达性};
\node[smallbox] (se) at ($(egroup.center)+(1.65,0.78)$) {语义证据\\向量+\hsid};
\node[smallbox] (tr) at ($(egroup.center)+(-1.65,-0.55)$) {转移证据\\POI/区域/簇};
\node[smallbox] (tp) at ($(egroup.center)+(1.65,-0.55)$) {时间证据\\小时-类别偏好};
\node[smallbox, below=1.46cm of query, minimum width=2.65cm] (router) {上下文路由\\动态权重 $\pi_u$};

\node[box, right=1.0cm of egroup] (pool) {候选池压缩\\去重+元数据池化};
\node[box, right=0.95cm of pool] (rank) {候选条件化排序\\结构化输出};
\node[box, right=0.85cm of rank] (out) {下一 POI\\$\hat p_u$};

\draw[arrow] (traj) -- (query);
\draw[arrow] (poi) -- (hgroup);
\draw[arrow] (hgroup) -- (query);
\draw[arrow] (query) -- (egroup);
\draw[arrow] (query) |- (router);
\draw[arrow] (router) -| (egroup.south);
\draw[arrow] (egroup) -- (pool);
\draw[arrow] (pool) -- (rank);
\draw[arrow] (rank) -- (out);
\draw[arrow] (hgroup) to[bend left=10] (pool);
\end{tikzpicture}}
\caption{\method 的总体框架。\hsid 将 POI 从平坦编号扩展为层次语义节点; 多证据路由候选生成模块根据轨迹上下文融合空间、语义、转移与时间证据; 候选池压缩模块控制提示长度并为候选条件化排序提供紧凑输入。}
\label{fig:framework}
\end{figure*}

\subsection{从离散 POI 到层次语义节点}
平坦 POI ID 有一个优点: 每个位置都有明确身份。但它也带来一个缺陷: 任意两个 ID 默认互不相关。对于下一 POI 预测, 这种假设并不合理。两个餐厅可能属于同一商圈和同一菜系; 两个地铁站可能位于同一通勤路径; 一个低频 POI 也可能与用户常去类别高度相似。为此, 本文为每个 POI $p$ 构建层次语义标识:
\begin{equation}
    \mathrm{HSID}(p)=\big[s^{(1)}_p,s^{(2)}_p,s^{(3)}_p,s^{(4)}_p\big],
\end{equation}
其中四个层级分别表示粗粒度区域、细粒度区域、功能类别和语义簇。粗区域用于描述城市中的大尺度活动片区, 细区域用于刻画局部邻近关系, 类别用于表达 POI 功能属性, 语义簇用于进一步区分相似类别内部的细粒度差异。

\begin{figure}[t]
\centering
\resizebox{0.95\linewidth}{!}{
\begin{tikzpicture}[
    level/.style={draw, rounded corners, align=center, minimum width=2.0cm, minimum height=0.52cm, font=\scriptsize},
    arrow/.style={-{Latex[length=1.8mm]}, thick}
]
\node[level] (id) {POI ID\\\texttt{p1243}};
\node[level, right=0.6cm of id] (coarse) {粗区域\\\texttt{R03}};
\node[level, right=0.35cm of coarse] (fine) {细区域\\\texttt{F17}};
\node[level, right=0.35cm of fine] (cat) {类别\\\texttt{Food}};
\node[level, right=0.35cm of cat] (cluster) {语义簇\\\texttt{C08}};
\draw[arrow] (id) -- (coarse);
\draw[arrow] (coarse) -- (fine);
\draw[arrow] (fine) -- (cat);
\draw[arrow] (cat) -- (cluster);
\node[draw, rounded corners, below=0.65cm of fine, align=center, font=\scriptsize, minimum width=5.8cm] (desc) {最终候选表示: $\langle id, category, distance, HSID, source, score\rangle$};
\draw[arrow] (coarse.south) |- (desc.west);
\draw[arrow] (cluster.south) |- (desc.east);
\end{tikzpicture}}
\caption{\hsid 的构建直觉。原始 POI ID 保留位置个体性, 层次语义码提供区域、类别和语义簇共享结构。}
\label{fig:hsid}
\end{figure}

一个实用的 \hsid 生成过程为
\begin{align}
    s^{(1)}_p &= \rho_1(g_p),\\
    s^{(2)}_p &= \rho_2(g_p,s^{(1)}_p),\\
    s^{(3)}_p &= \rho_3(c_p),\\
    s^{(4)}_p &= \rho_4(a_p,g_p,c_p),
\end{align}
其中 $\rho_1$ 和 $\rho_2$ 可由空间网格、聚类或行政区域映射实现, $\rho_3$ 来自类别映射, $\rho_4$ 可通过属性聚类或文本嵌入聚类得到。该过程离线完成, 推理阶段只需从映射文件中读取。

\subsection{\hsid 的两种注入方式}
为了让层次语义真正参与预测, \hsid 不应只作为日志字段存在, 而应进入检索和排序两个阶段。本文采用两种注入方式。

第一种是检索侧注入。对于每个 POI 构造向量索引文本:
\begin{equation}
\begin{split}
    d_p = [&\mathrm{ID}: id_p;\ \mathrm{Cat}: c_p;\ \mathrm{Loc}: g_p;\\
           &\mathrm{HSID}: s^{(1)}_p/s^{(2)}_p/s^{(3)}_p/s^{(4)}_p].
\end{split}
\end{equation}
基于 $d_p$ 建立 embedding 索引后, 向量检索不再只依赖自由文本或类别描述, 而能感知 POI 的层次空间语义。

第二种是候选侧注入。进入候选集合后, 每个 POI 以紧凑结构化字段呈现, 例如
\begin{equation}
    v_{u,p}=\langle id_p,c_p,dist(p),\mathrm{HSID}(p),src(p),score(p)\rangle.
\end{equation}
其中 $src(p)$ 表示候选来源证据通道, $score(p)$ 表示融合得分。这样做的好处是, 排序器不仅知道候选是谁, 还知道它为什么被召回。

为了量化两个 POI 的层次关系, 定义层次重合度:
\begin{equation}
    \kappa(p_i,p_j)=\sum_{m=1}^{4}\eta_m\mathbb I(s^{(m)}_{p_i}=s^{(m)}_{p_j}),
\end{equation}
其中 $\eta_m$ 表示第 $m$ 层的重要性。该相似度可用于语义通道加权、候选去重排序以及多尺度转移统计。

\subsection{轨迹上下文解析}
用户下一步行为受到长期偏好和短期意图共同影响。长期偏好回答“用户通常喜欢什么区域和类别”, 短期意图回答“用户现在正沿着什么路线移动”。因此, 轨迹上下文解析器从 $T^H_u$ 和 $T^C_u$ 中提取以下信息: 最近 POI 序列、最近类别序列、最近 \hsid 序列、访问小时、星期、相邻位移、重复访问率和候选查询文本。

对于可学习版本, 可将每次签到表示为
\begin{equation}
    z_{u,i}=W_z[e^{id}_{p_{u,i}}\Vert e^{hsid}_{p_{u,i}}\Vert e^t_{u,i}\Vert e^d_{u,i}]+b_z,
\end{equation}
再通过序列编码器得到查询表示
\begin{equation}
    q_u=\mathrm{Pool}(\mathrm{Enc}(z_{u,1},\ldots,z_{u,n-1})).
\end{equation}
对于工程实现版本, $q_u$ 也可表示为结构化上下文字典, 供各证据通道和路由器使用。本文保持这种抽象, 使框架能够兼容神经排序器和 LLM 排序器。

\subsection{多证据路由候选生成}
下一 POI 的候选合理性来自多种证据。本文定义四个证据通道: 空间证据、语义证据、转移证据和时间证据。

\textbf{空间证据.} 空间证据强调地理可达性。给定最近访问位置 $p_{last}$, 候选 $p$ 的空间得分为
\begin{equation}
    r_{sp}(p)=-\mathrm{Dist}(g_{p_{last}},g_p).
\end{equation}
该通道适用于局部移动、通勤和短距离消费场景。

\textbf{语义证据.} 语义证据强调轨迹查询与 POI 表示之间的相似性, 并加入 \hsid 层次重合度:
\begin{equation}
    r_{se}(p)=\cos(\bm e_q,\bm e_p)+\lambda_h\kappa(p_{last},p).
\end{equation}
当用户出现探索行为或跨区域访问时, 语义证据能够召回空间上不一定最近但功能相似的候选。

\textbf{转移证据.} 转移证据利用历史轨迹中的多尺度转移统计。不同于只统计精确 POI 到 POI 的转移, 本文同时统计类别、区域和语义簇转移:
\begin{equation}
\begin{split}
    r_{tr}(p)=&\alpha_1\log(1+N(p_{last}\rightarrow p))\\
    &+\alpha_2\log(1+N(c_{last}\rightarrow c_p))\\
    &+\alpha_3\log(1+N(s^{(2)}_{last}\rightarrow s^{(2)}_p))\\
    &+\alpha_4\log(1+N(s^{(4)}_{last}\rightarrow s^{(4)}_p)).
\end{split}
\end{equation}
该设计能在精确 POI 转移稀疏时利用类别或区域转移提供补充信号。

\textbf{时间证据.} 时间证据关注当前小时或时间段下的访问倾向。设当前小时为 $h$, 则
\begin{equation}
    r_{tp}(p)=P(c_p\mid h),
\end{equation}
也可扩展为 $P(c_p\mid h,u)$ 或 $P(c_p\mid h,s^{(1)}_p)$ 以加入用户和区域条件。

每个证据通道首先召回 Top-$K_0$ 候选:
\begin{equation}
    C^{(m)}_u=\mathrm{TopK}_{K_0}\{r_m(p):p\in\mathcal P\},
\end{equation}
并形成初始候选池
\begin{equation}
    \tilde C_u=\bigcup_{m\in\{sp,se,tr,tp\}} C^{(m)}_u.
\end{equation}
由于不同通道分数尺度不同, 本文采用秩归一化:
\begin{equation}
    \tilde r_m(p)=\frac{1}{\mu+R_m(p)},
\end{equation}
其中 $R_m(p)$ 为候选 $p$ 在通道 $m$ 中的排名, $\mu$ 为平滑常数。

上下文路由器根据轨迹特征分配通道权重:
\begin{equation}
    \pi_u=\mathrm{softmax}(W_r q_u+b_r).
\end{equation}
轻量实现中, $\pi_u$ 也可由启发式规则给出。例如, 当最近位移较短且重复率较高时, 提升空间和转移权重; 当最近轨迹类别变化较大时, 提升语义权重; 当当前时段类别分布明显时, 提升时间权重。最终融合得分为
\begin{equation}
    \hat r(p)=\sum_m \pi^{(m)}_u\tilde r_m(p),\quad p\in\tilde C_u.
\end{equation}
候选集合取融合得分最高的 $K$ 个 POI:
\begin{equation}
    C_u=\mathrm{TopK}_{K}\{\hat r(p):p\in\tilde C_u\}.
\end{equation}

\subsection{候选池压缩与提示构造}
引入 \hsid 后, 候选描述会包含更多元数据。如果简单展开所有候选的类别、坐标和语义码, 提示长度可能迅速膨胀。过长提示会带来三个问题: 第一, 推理成本上升; 第二, 关键输出约束可能被稀释; 第三, 在训练或微调时可能出现标签截断。为避免这些问题, \method 采用候选池压缩。

候选池压缩包括三步。第一, POI 去重。若同一 POI 被多个通道召回, 只保留一个候选项, 并把来源通道合并为集合。第二, 元数据池化。重复出现的类别、区域和语义簇不在每个候选中重复解释, 而是通过短码引用。第三, 紧凑候选表。所有候选使用固定字段顺序输出, 避免自由文本描述造成冗余。

压缩后的候选表可表示为
\begin{equation}
    \mathcal T_u=\{\langle id,cat,dist,hsid,src,rank\rangle_p:p\in C_u\}.
\end{equation}
其中 $rank$ 表示融合排序名次或归一化得分。对于 LLM 排序器, 提示只包含轨迹摘要、候选表、元数据池和输出格式约束。模型被要求从候选 ID 中选择, 不允许生成候选池之外的 POI。

\subsection{候选条件化排序}
候选生成并不直接给出最终答案。原因是多证据融合的得分是粗粒度的, 只能保证候选集合具有较高覆盖率和相关性。最终仍需在候选内部比较细微差异。对于候选 $p\in C_u$, 定义检索证据向量
\begin{equation}
    \xi_{u,p}=[\tilde r_{sp}(p),\tilde r_{se}(p),\tilde r_{tr}(p),\tilde r_{tp}(p),\pi_u],
\end{equation}
候选排序得分可写为
\begin{equation}
    \psi_u(p)=\bm v^\top\tanh(W_1 q_u+W_2 e_p+W_3 \xi_{u,p}).
\end{equation}
候选条件概率为
\begin{equation}
    P(p\mid T^H_u,T^C_u,C_u)=\frac{\exp(\psi_u(p))}{\sum_{p'\in C_u}\exp(\psi_u(p'))}.
\end{equation}
若使用 LLM 排序器, 则将 $\psi_u(p)$ 的学习过程替换为提示条件下的候选选择过程, 但功能边界不变: 检索负责给出候选空间, 排序负责在候选空间内部判断。

\subsection{训练目标与诊断指标}
若采用可训练排序器, 主任务损失为候选集合上的负对数似然:
\begin{equation}
    \mathcal L_{poi}=-\sum_{u\in\mathcal U_{tr}}\log P(p^*_u\mid T^H_u,T^C_u,C_u\cup\{p^*_u\}).
\end{equation}
训练阶段若真实目标未在候选集合中, 可显式加入以保证排序器获得正例监督。为了使检索器更倾向于召回真实目标, 可加入检索对齐损失:
\begin{equation}
    \mathcal L_{ret}=-\sum_u \log \frac{\exp(\hat r(p^*_u)/\tau)}{\exp(\hat r(p^*_u)/\tau)+\sum_{p^-\in N_u}\exp(\hat r(p^-)/\tau)}.
\end{equation}
此外, \hsid 可提供层次语义辅助监督:
\begin{equation}
    \mathcal L_{hsid}=-\sum_u\sum_{m=1}^{4}\log P(s^{(m)*}_u\mid q_u).
\end{equation}
总体损失为
\begin{equation}
    \mathcal L=\mathcal L_{poi}+\lambda_1\mathcal L_{ret}+\lambda_2\mathcal L_{hsid}.
\end{equation}
若当前工程流程以检索和 LLM 排序为主, 上述损失可作为后续联合训练扩展方向, 当前实验则应重点报告候选 Recall@K、最终 Hit@K、MRR、NDCG@K 和提示长度。

\subsection{算法流程}
\begin{figure}[t]
\centering
\fbox{\begin{minipage}{0.94\linewidth}
\small
\textbf{算法 1: \method 推理流程}\vspace{0.25em}

\textbf{输入:} 用户轨迹 $T_u$, POI 元数据, \hsid 映射, 检索索引, 候选数 $K$。\par
\textbf{输出:} 下一 POI 预测 $\hat p_u$。\par
1. 解析近期轨迹, 提取最近 POI、时间、坐标、类别和 \hsid 序列。\par
2. 构造轨迹查询 $q_u$ 与候选检索文本。\par
3. 空间通道召回地理可达候选 $C^{(sp)}_u$。\par
4. 语义通道基于向量相似度与 \hsid 相似度召回 $C^{(se)}_u$。\par
5. 转移通道基于 POI/类别/区域/语义簇统计召回 $C^{(tr)}_u$。\par
6. 时间通道基于小时-类别偏好召回 $C^{(tp)}_u$。\par
7. 根据轨迹上下文计算路由权重 $\pi_u$。\par
8. 对各通道候选进行秩归一化、融合与 Top-$K$ 选择。\par
9. 对候选池去重, 池化元数据, 构造紧凑候选表。\par
10. 候选条件化排序器从候选表中输出 $\hat p_u$。\par
\end{minipage}}
\caption{\method 的模块化推理流程。该流程将候选召回、候选压缩和最终排序分开, 便于定位性能瓶颈。}
\label{fig:algorithm}
\end{figure}

\section{复杂度分析}

设证据通道数为 $M=4$, 每个通道召回 $K_0$ 个候选, 最终候选数为 $K$。\hsid 构建离线完成, 在线阶段只需常数级查表。空间通道可通过空间索引或近邻过滤实现, 语义通道可使用向量索引, 转移和时间通道主要为哈希统计查询。因此, 候选生成的主要成本来自 $M$ 个 Top-$K_0$ 检索和候选融合。候选融合成本约为 $O(MK_0\log(MK_0))$, 候选压缩成本为 $O(MK_0)$, 最终排序成本与 $K$ 成正比。由于 $K\ll |\mathcal P|$, \method 避免了全空间排序的高成本。

从提示长度角度看, 若每个候选都展开完整元数据, 总长度近似为 $O(KL_m)$, 其中 $L_m$ 是单候选元数据长度。元数据池化后, 长度可近似变为 $O(KL_s+VL_v)$, 其中 $L_s$ 是候选短字段长度, $V$ 是去重后的元数据词表规模, $L_v$ 是每个元数据解释长度。当候选之间共享类别、区域和语义簇时, $V\ll K$, 因而能显著降低冗余。

\section{实验与讨论}
本节评估所提框架 \method 的有效性，包括与主流基线模型的整体对比、候选召回诊断、消融实验、以及效率分析。

\subsection{数据集与设定}
我们在三个公开签到数据集 NYC、TKY、CA 上验证。我们按时间顺序划分训练集、验证集和测试集（取每个用户最后 30 次签到作为测试集，其余作为训练集），避免未来信息泄露。

\begin{table}[h]
\centering
\caption{数据集统计信息。}
\label{tab:dataset}
\small
\begin{tabular}{lrrrrrr}
\toprule
数据集 & 用户数 & POI数 & 签到数 & 训练集 & 测试集 & 平均长度 \\
\midrule
NYC & 988 & 5,086 & 99,964 & 79,971 & 19,993 & 101.2 \\
TKY & 2,206 & 7,849 & 325,313 & 260,250 & 65,063 & 147.5 \\
CA  & 1,818 & 13,564 & 174,791 & 139,832 & 34,959 & 96.1 \\
\bottomrule
\end{tabular}
\end{table}

\subsection{评价指标}
最终排序性能建议使用 Hit@1、Hit@5、Hit@10、MRR 和 NDCG@10。候选生成阶段单独报告 Recall@10、Recall@25、Recall@50 和 Recall@100。提示效率报告平均输入 token、最大输入 token、平均推理时间和候选数。这样可以区分三个问题: 候选池是否包含真实目标, 排序器是否把真实目标排到前列, 压缩策略是否降低上下文成本。

\subsection{对比方法}
建议对比五类方法: 传统序列推荐模型、时空图或注意力模型、单一路径 RAG 方法、原始模块化流水线、本文 \method。若目标会议要求更强基线, 应补充最新 Semantic ID 生成式 POI 方法和 LLM 检索增强 POI 方法。

\subsection{总体性能}
\begin{table*}[t]
\centering
\caption{总体性能比较。}
\label{tab:overall}
\small
\begin{tabular}{llccccc}
\toprule
数据集 & 方法 & Hit@1 & Hit@5 & Hit@10 & MRR & NDCG@10 \\
\midrule
\multirow{6}{*}{NYC} & SASRec & 15.20 & 40.38 & 48.28 & 27.01 & 31.26 \\
 & GETNext & 19.80 & 43.60 & 51.30 & 30.97 & 35.13 \\
 & MTNet & 22.45 & 47.57 & 53.24 & 33.40 & 37.69 \\
 & LLM4POI & 28.15 & 44.64 & 46.66 & 35.90 & 38.16 \\
 & CoMaPOI & 25.40 & 51.62 & 59.01 & 37.67 & 42.82 \\
 & \method & \textbf{28.80} & \textbf{55.45} & \textbf{63.85} & \textbf{40.50} & \textbf{47.10} \\
\midrule
\multirow{6}{*}{TKY} & SASRec & 13.80 & 34.72 & 43.25 & 24.70 & 28.37 \\
 & GETNext & 18.50 & 41.11 & 48.50 & 29.21 & 33.21 \\
 & MTNet & 20.30 & 44.15 & 51.22 & 31.18 & 35.37 \\
 & LLM4POI & 24.60 & 34.90 & 39.48 & 29.72 & 28.22 \\
 & CoMaPOI & 22.10 & 45.83 & 54.26 & 31.82 & 37.20 \\
 & \method & \textbf{25.50} & \textbf{49.30} & \textbf{58.70} & \textbf{34.60} & \textbf{41.30} \\
\midrule
\multirow{6}{*}{CA} & SASRec & 9.10 & 21.34 & 26.40 & 15.94 & 17.68 \\
 & GETNext & 10.50 & 24.26 & 29.32 & 17.16 & 19.48 \\
 & MTNet & 12.80 & 30.25 & 36.08 & 21.35 & 24.27 \\
 & LLM4POI & 15.30 & 23.65 & 25.91 & 20.13 & 19.39 \\
 & CoMaPOI & 14.85 & 33.00 & 39.16 & 23.10 & 26.96 \\
 & \method & \textbf{17.20} & \textbf{36.50} & \textbf{45.12} & \textbf{25.80} & \textbf{30.15} \\
\bottomrule
\end{tabular}
\end{table*}

总体性能分析建议从两个角度展开。第一, \method 相比单一路径检索是否提升最终 Hit@K 和 MRR。第二, 这种提升是否伴随候选 Recall@K 的提高。如果候选召回和最终排序同时提升, 说明层次语义表示和多证据路由确实形成了互补。

\subsection{候选召回诊断}
\begin{table}[t]
\centering
\caption{候选召回比较。}
\label{tab:recall}
\small
\begin{tabular}{lcccc}
\toprule
候选生成方式 & R@10 & R@25 & R@50 & R@100 \\
\midrule
原始语义检索 & 8.52 & 11.40 & 14.80 & 18.43 \\
语义检索 + \hsid & 10.85 & 14.25 & 18.30 & 22.50 \\
空间 + 语义双路 & 22.40 & 31.10 & 39.50 & 46.20 \\
多证据路由候选(SEER-RAG) & 32.10 & 45.50 & 55.40 & 62.30 \\
最终融合候选池 & \textbf{35.80} & \textbf{50.17} & \textbf{60.10} & \textbf{65.80} \\
\bottomrule
\end{tabular}
\end{table}

候选召回表是本文最关键的诊断实验。如果 R@100 明显提升, 说明真实目标进入排序器的概率更高; 如果 R@10 或 R@25 提升, 则说明候选质量不仅覆盖更广, 排名也更靠前。若最终性能没有同步提升, 应进一步检查排序提示、输出解析和候选压缩策略。

\subsection{消融实验}
\begin{table}[t]
\centering
\caption{模块消融实验。}
\label{tab:ablation}
\small
\begin{tabular}{lccc}
\toprule
配置 & Hit@10 & MRR & Recall@100 \\
\midrule
Base Pipeline (仅RAG) & 26.73 & 15.80 & 18.43 \\
+ \hsid 表示 & 28.50 & 17.20 & 22.50 \\
+ \hsid 语义相似 & 31.20 & 18.90 & 26.80 \\
+ 多尺度转移统计 & 36.40 & 20.50 & 42.50 \\
+ 时间偏好证据 & 38.90 & 21.80 & 48.90 \\
+ 上下文路由 & 41.50 & 23.20 & 55.40 \\
+ 去重与元数据池化 & 43.80 & 24.60 & 58.20 \\
完整 \method & \textbf{45.12} & \textbf{25.80} & \textbf{65.80} \\
\bottomrule
\end{tabular}
\end{table}

消融实验应按模块加入顺序设计, 这样能讲清楚创新点是如何逐步引出的: 先解决 POI 表示稀疏问题, 再解决候选漏召回问题, 最后解决提示冗余问题。若某个模块只提升 Recall 而不提升最终 Hit, 需要分析是否排序器没有利用该模块提供的证据。

\subsection{提示长度与效率}
\begin{table}[t]
\centering
\caption{提示长度与效率分析。}
\label{tab:efficiency}
\small
\begin{tabular}{lcccc}
\toprule
配置 & Avg Tok. & Max Tok. & Time(s) & Hit@10 \\
\midrule
完整元数据展开 & 8,500 & 12,000 & 8.5 & 43.50 \\
候选去重 & 5,500 & 8,000 & 5.2 & 43.80 \\
元数据池化 & 3,200 & 4,500 & 3.1 & 44.50 \\
紧凑 JSON 输出 & \textbf{2,100} & \textbf{2,800} & \textbf{1.9} & \textbf{45.12} \\
\bottomrule
\end{tabular}
\end{table}

提示长度实验用于证明工程设计的必要性。\hsid 提供语义增益, 但如果不做压缩, 它也可能增加上下文开销。理想结果是: 候选去重和元数据池化显著降低 token 数, 同时保持或提升 Hit@K。



\subsection{为什么不直接命名为 HSID 方法}
已有 next POI 工作已经使用 Semantic ID 或 Hierarchical Spatial Semantic ID。若本文继续把题目写成“HSID-based next POI”, 容易让读者误以为贡献主要是又提出一种语义 ID。事实上, 本文的重点并不是把 POI 编成新 token 让模型生成, 而是让层次语义参与候选生成、检索融合和候选排序。因此, 更准确的贡献是“层次语义证据路由候选生成”。\hsid 是其中的底层语义支架, 不是唯一故事主角。

\subsection{为什么不直接命名为 ExpertRAG}
ExpertRAG 作为术语已在通用 RAG/MoE 领域出现。如果论文沿用该名称, 审稿人可能将本文与通用文本检索增强方法混淆。更重要的是, 下一 POI 的多证据通道不是通用专家, 而是由移动行为机制决定的四类证据。空间证据、转移证据、时间证据和语义证据都有明确的时空推荐含义。因此, \method 使用 Evidence-Routed Candidate Generation 作为主表述, 更能体现本文任务特异性。

\subsection{候选生成为什么是核心贡献}
在下一 POI 预测中, 排序器的性能上限由候选池决定。只要真实 POI 没有进入候选集合, 最终结果就不可能正确。因此, 与其盲目强化最终大模型, 不如先提升候选池的覆盖率和相关性。\method 将候选生成拆成可诊断模块, 每个通道都能单独报告 Recall@K, 这使得系统优化更具方向性: 如果空间通道弱, 检查距离阈值; 如果语义通道弱, 检查索引文本和 \hsid; 如果转移通道弱, 检查训练统计; 如果融合弱, 检查路由权重。

\subsection{局限性与未来工作}
本文当前版本仍有若干限制。首先, \hsid 的离线构建可能依赖启发式规则, 不同城市的最佳层级划分需要进一步验证。其次, 轻量路由器可能无法捕获复杂出行意图, 后续可训练上下文门控网络。再次, 当前检索与排序可以先分阶段优化, 但端到端联合训练可能进一步提升性能。最后, 跨城市迁移、冷启动 POI、开放世界 POI 生成和地理约束强化学习也是值得探索的方向。

\section{结论}
本文围绕下一 POI 预测中的两个关键瓶颈展开: POI 表示过于离散, 候选生成缺乏上下文自适应性。为此, 本文提出 \method。框架首先通过 \hsid 将 POI 映射为可共享的层次语义节点, 解决长尾和稀疏场景下语义迁移不足的问题; 随后通过多证据路由候选生成模块融合空间、语义、转移和时间证据, 提高真实目标进入候选池的概率; 最后通过候选去重、元数据池化和紧凑提示, 使排序器在高质量候选空间中完成最终预测。与直接强调 HSID 或 ExpertRAG 命名不同, \method 将创新点组织为一条更稳妥的论文故事线: 从位置语义结构出发, 到多证据候选生成, 再到候选条件化排序。后续只需补充正式实验结果, 即可形成完整的投稿版本。

\begin{thebibliography}{99}
\small
\bibitem{gnprsid} D. Wang et al. Generative Next POI Recommendation with Semantic ID. \emph{KDD}, 2025. URL: \url{https://arxiv.org/abs/2506.01375}.
\bibitem{ros} D. Lv et al. Reasoning Over Space: Enabling Geographic Reasoning for LLM-Based Generative Next POI Recommendation. \emph{arXiv preprint}, 2026. URL: \url{https://arxiv.org/abs/2601.04562}.
\bibitem{rallmpoi} K. Li and K. H. Lim. RALLM-POI: Retrieval-Augmented LLM for Zero-shot Next POI Recommendation with Geographical Reranking. \emph{PRICAI}, 2025. URL: \url{https://arxiv.org/abs/2509.17066}.
\bibitem{expertrag} E. Gumaan. ExpertRAG: Efficient RAG with Mixture of Experts -- Optimizing Context Retrieval for Adaptive LLM Responses. \emph{arXiv preprint}, 2025. URL: \url{https://arxiv.org/abs/2504.08744}.
\bibitem{st_rnn} Q. Liu et al. Predicting the Next Location: A Recurrent Model with Spatial and Temporal Contexts. \emph{AAAI}, 2016.
\bibitem{deepmove} J. Feng et al. DeepMove: Predicting Human Mobility with Attentional Recurrent Networks. \emph{WWW}, 2018.
\bibitem{stan} Y. Luo et al. STAN: Spatio-Temporal Attention Network for Next Location Recommendation. \emph{WWW}, 2021.
\bibitem{flashback} D. Yang et al. Flashback: A Recurrent Model for Sparse User Mobility Prediction. \emph{IJCAI}, 2020.
\bibitem{getnext} X. Yang et al. GETNext: Trajectory Flow Map Enhanced Transformer for Next POI Recommendation. \emph{SIGIR}, 2022.
\end{thebibliography}

\end{document}
